\SourceLeaf{482}
\Needspace{6\baselineskip}\begin{center}\bfseries SCHLUSSBEMERKUNG ZUR NEUEN AUSGABE.\end{center}

Dieser erste Band von \textsc{Gauss} Werken ist ein Wiederabdruck der im Jahre 1801 in Octav erschienenen sieben Sectionen der Disquis. Arithm. Die achte Section, auf die an mehreren Stellen verwiesen wird und die \textsc{Gauss} wol anfangs mit den übrigen zu veröffentlichen beabsichtigte, findet sich unter seinen Handschriften. Da er die Ausarbeitung derselben aber nicht in gleicher Weise abgeschlossen hat wie die der sieben ersten Sectionen, so wird sie in dieser Ausgabe den arithmetischen Abhandlungen des Nachlasses sich anschliessen.

Textänderungen sind in den Disqu. Ar. nur an folgenden Stellen vorgenommen:

In Art. 125 sind die beiden Einschaltungen (si \(>5\)) und (\(>17\); sed \(-13N3,\allowbreak\ -17N5\)) eingefügt worden.

In Art. 126. \textit{Demonstr.} ist ‘in serie (I) \(a\) terminos esse per \(p\) divisibiles, \(b\) terminos per \(p^2\) divisibiles, \(c\) terminos per \(p^3\) divisibiles etc.’ statt ‘in serie (I) \(a\) terminos esse per \(p\) divisibiles neque vero per \(p^2\), \(b\) terminos per \(p^2\) non autem per \(p^3\) divisibiles, \(c\) terminos per \(p^3\) non autem per \(p^4\) etc.’ gesetzt worden.

In Art. 128 III sind die nach ‘Si enim \(+b\) vel \(-b\equiv r\ (\mathrm{mod.}\ p)\), erit \(bb\equiv rr\ (\mathrm{mod.}\ 2p)\), adeoque terminus \(\frac12(a-bb)\) per \(p\) divisibilis.’ in der Ausgabe von 1801 noch folgenden Worte ‘multoque magis \(2(a-bb)\).’ ausgelassen.

In Art. 129 ist überall \(2\surd a+1\) statt \(2\surd a\) und demnach 9 statt 4 als diejenige Zahl, bis zu welcher \(a\) jedenfalls nicht herabgeht, gesetzt worden.

In Art. 139 lautete der Schluss der Untersuchung über den Fall (4) ‘Facile vero perspicitur, ex ista aequatione deduci posse haec \(a'pRh,\allowbreak\ \pm ahpRa',\allowbreak\ \pm aa'hRp\); quae cum iis quae in (2) invenimus conveniunt. In reliquis autem demonstratio est eadem.’ Dieses ist gemäss der Note geändert, die \textsc{Gauss} dem Art. 2 der Abhandlung \textit{Theorematis arithmetici demonstratio nova} beigefügt hat: ‘Haud abs re erit, levem aliquem errorem, qui nescio qua negligentia in illius demonstrationis expositionem irrepsit, hic indicare atque corrigere. Pag. (108) inde a l. (14) ratiocinia sequentia sunt substituenda: Facile vero perspicitur, ex ista aequatione deduci posse haec \(a'pRh\ldots(\alpha);\allowbreak\ \pm ahRa'\ldots(\delta);\allowbreak\ \pm ahRp\ldots(\gamma)\). Ex \((\alpha)\) quod convenit cum \((\alpha)\) in (2) sequitur perinde ut illic, esse vel simul \(hRp,\allowbreak\ hRa'\), vel \(hNp,\allowbreak\ hNa'\). Sed in casu priori foret per \((\delta)\), \(aRa'\) contra hyp.; quare erit \(hNp\), adeoque per \((\gamma)\) etiam \(aNp\).

\SourceLeaf{483}
In Art. 171 ist zufolge einer handschriftlichen Bemerkung ‘in qua \(A\) nec maior quam \(\surd{\frac43D},\allowbreak\ C\), nec minor quam \(2B\)’ statt ‘in qua \(A\) non \(>\surd{\frac43D}\), \(B\) non \(>\frac12A\), \(C\) non \(<A\)’ gesetzt worden.

In Art. 174. \textit{Methodus secunda.} ist ‘Quum \(a\) non erit \(>\surd{\frac43D}\), omnes formae quae hoc modo prodeunt, manifesto erunt reductae.’ gesetzt statt ‘Si quae formae hoc modo prodeunt, in quibus \(a>\surd{\frac43D}\), erunt reiiciendae, reliquae vero omnes manifesto erunt reductae.’

In Art. 256 VI enthält die frühere Ausgabe unter den positiven Determinanten von der Form \(8n+5\), für welche die Anzahl der Classen in der eigentlich primitiven Ordnung dreimal grösser ist als die in der uneigentlich primitiven, auch den Determinant 397, dem aber von beiden Arten gleich viel Classen zugehören.

In Art. 302 ist zufolge einer handschriftlichen Bemerkung die Anzahl der Classen für das zweite Tausend der negativen Determinanten zu 28595 angegeben und nicht zu 28603, wie in der frühern Ausgabe steht.

In Art. 303 sind für die Classificationen II. 4; II. 5; IV. 2 die Mengen der Determinanten nach den im Nachlass vorgefundenen Tafeln zu 31, 44, 69 angegeben und nicht wie in der früheren Ausgabe zu 32, 42, 68.

In Art. 325 ist ‘\(23t,\allowbreak\ 23t\pm1,\allowbreak\ 23t\pm5,\allowbreak\ 23t\pm7,\allowbreak\ 23t\pm9,\allowbreak\ 23t\pm10\). His deletis superstites inveniuntur 119, 127, qui duo soli ipsi \(V\) valorem quadratum conciliant’, gesetzt statt ‘\(23t,\allowbreak\ 23t\pm5,\allowbreak\ 23t\pm7\), \(23t\pm9,\allowbreak\ 23t\pm10\). His deletis superstites inveniuntur 119, 127, 137, e quibus duo priores soli ipsi \(V\) valorem quadratum conciliant.’

In Art. 360 IV sind die Worte ‘quum pro plerisque aliis aequationibus cubicis, quarum radices omnes reales sunt, simul anguli et rationis trisectio evitari nequeat.’ die auf: ‘e. g. pro \(\beta=3\) sola trisectio anguli’, folgten, und deren Unrichtigkeit \textsc{Gauss} in seinem Handexemplare notirt hat, ausgelassen.

Die Noten auf Seite 80, 102, 133, 211, 263, 265, 449, sind dem handschriftlichen Nachlasse entlehnt.

Die äussere Form ist bei der neuen Ausgabe zur Erleichterung der Uebersicht einigen Abänderungen gegen den Druck dieses Werkes im Jahre 1801 unterworfen, man glaubte sich dazu um so mehr berechtigt, als \textsc{Gauss} ausgesprochener Weise auch in anderen Punkten auf Raumersparniss Rücksicht genommen. Viele Formeln, die der Textzeile eingeschlossen waren, sind abgesondert herausgesetzt. Die Inhaltsangaben, die zusammengestellt dem Werke vorangingen, sind auch den einzelnen Abtheilungen und nicht nur den Sectionen wie in der ersten Ausgabe beigefügt, die allgemeineren darunter sind den Seiten als Überschrift gegeben.

\BookRule{45mm}

\begin{center}\small GÖTTINGEN,\\ GEDRUCKT IN DER DIETERICHSCHEN UNIVERSITÄTS-DRUCKEREI\\ W. FR. KAESTNER.\end{center}
